% ----------------------
\subsection{Section 118 (8 July 1838)}
% ----------------------
	\label{DC118}
	
	% -----
	\subsubsection{Historical Background}
	% -----
		\label{DC118-HB}
		
		% --
		\paragraph{JSP}
		% --
		
			``On Sunday, 8 July 1838, JS dictated five revelations in Far West, Missouri, each of which concerned church leadership or finances; one of these revelations regarded the Quorum of the Twelve Apostles. The quorum had not been immune to the dissent and disaffiliation that plagued the church in 1837 and 1838. Apostles Luke Johnson and John F. Boynton renounced the church in December 1837 and were consequently excommunicated by the high council. In the quarterly conference held 7--8 April 1838, David W. Patten gave a report on the quorum and stated that he could not recommend Lyman Johnson or William E. McLellin. Lyman Johnson, who may have sympathized with his brother Luke, was excommunicated for various transgressions on 13 April 1838. McLellin, who had been troubled with the church and JS for some time, was `found in transgression' in a church trial held 11 May. JS apparently began selecting replacements for disaffected members of the quorum before he departed Kirtland, Ohio, for Far West on 12 January 1838. For example, John E. Page and John Taylor had been designated as replacements for Luke Johnson and Boynton by early January 1838.
			
			``Despite the turmoil in the church, most of the apostles remained loyal to JS or repaired their relationships with him, and many were serving missions or were expecting to serve in the near future. Under JS's direction, apostles Heber C. Kimball and Orson Hyde undertook a dramatically successful mission to England from June 1837 to May 1838. During the year prior to the dictation of this 8 July 1838 revelation, JS dictated revelations for senior apostles Thomas B. Marsh, David W. Patten, and Brigham Young, with each revelation mentioning or implying forthcoming proselytizing assignments. The revelation for Patten directed him to prepare to embark on a mission the following spring and implied that the other apostles would go with him. On 6 July 1838, JS received a letter from Kimball and Hyde reporting on their return from England and the hundreds of new converts they had brought into the church.
			
			``Two days later, in a leadership meeting held before the Sunday worship services for `the congregation of the saints,' JS dictated this revelation regarding the Quorum of the Twelve Apostles. The Sunday leadership meeting was probably associated with the quarterly conference held the previous two days. When Thomas B. Marsh wrote about the leadership meeting to Wilford Woodruff, he recounted that `Prest. Joseph Smith J\uline{r.} and somome others ware assembled togeather to attend to some church business.' In JS's journal, George W. Robinson wrote that the revelation was `given to the Twelve Apostles'; however, Marsh was the only apostle present at the meeting. According to JS's journal, the meeting included `J smith Jr. S[idney] Rigdon, H[yrum] smith, E[dward] Partridge I, Morly [Isaac Morley] J[ared] Carter, S[ampson] Avard T[homas] B, Marsh \& G[eorge] W, Robinson,' all of whom held positions of leadership. JS, Rigdon and Hyrum Smith composed the First Presidency. Partridge was the bishop of Zion, and Morley was the first counselor in the bishopric. Carter was the captain general of the Society of the Daughter of Zion (Danites), and Avard was the society's brigadier general. Marsh was the pro tempore president of Zion and the president of the Quorum of the Twelve Apostles. Robinson, the church recorder and clerk as well as the scribe for the First Presidency, was probably present in a clerical capacity.
			
			``In JS's journal, this revelation appears as the first of five revelations dictated on 8 July. JS's journal entry for 8 July introduces the first three revelations by stating they were read to the congregation that met later that day but does not specify whether the fourth and fifth revelations were read at the same time. The fifth revelation---addressed to William Marks, Newel K. Whitney, and Oliver Granger---was copied into a letter the First Presidency wrote later in the day to Marks and Whitney. The letter states that the revelation had been `rec\supersu{d.} this morning.' The order in which the revelations were copied into JS's journal suggests that all five revelations were dictated on the morning of 8 July, apparently in the leadership meeting mentioned in the introduction to the first revelation.
			
			``In the letter that Marsh wrote to Woodruff, he explained that JS dictated this revelation `when it was thought proper to select those who ware designed of the Lord to fill the places of those of the twelve who had fallen.' As presented in JS's journal, the revelation was a direct response to the plea to `show unto us thy will O, Lord concerning the Twelve.' The revelation named new apostles---Taylor, Page, Woodruff, and Willard Richards---to replace those who had been removed, and it also directed all the members of the quorum to prepare for a mission `over the great waters' the following spring. This and the other 8 July revelations were probably transcribed by Robinson as JS dictated them.
			
			``The text of this 8 July revelation was read aloud in the worship meeting for all church members later in the day. Marsh responded to the revelation by calling for a meeting the next day with Sidney Rigdon and the four other members of the quorum who were available in the area. They agreed to contact absent and newly appointed members of the quorum to inform them of their expected mission abroad.
			
			``In addition to the copy of the revelation Robinson transcribed in JS's journal, copies were also made by Brigham Young, Wilford Woodruff, and Willard Richards. A comparison of the copies suggests that Young's version most closely represents the wording of the original revelation; therefore, this version is featured here.%
				\footnote{%
					% \label{}%
					JSP note: ``Comparisons of the versions of the five 8 July 1838 revelations show that Robinson added introductory phrases and made other slight revisions to polish the texts when he copied them into JS's journal. ''
					}
			Young inscribed the revelation in a book that he intermittently used as a notebook and a journal. Contextual evidence suggests that Young copied the revelation between April 1839 and 12 October 1840. He apparently copied two other documents at the time that he copied the revelation. One of these documents, a 30 March 1839 letter Orson Hyde wrote from Missouri to Young in Illinois indicates Young did not copy the revelation until at least April 1839, after receiving the letter from Hyde. Immediately following the three copied documents, Richards inscribed a page and a half of his own genealogical information. Just below this information, Young recorded a journal entry for 12 October 1840. Because the genealogical information written by Richards included dates from 1839, Richards apparently inscribed it sometime after Young joined him in England in April 1840 but no later than 12 October 1840, when Young added the new journal entry. Based on the reception of Hyde's letter and the inscription of Richards's genealogical information, it can be determined that Young copied the three documents sometime between April 1839 and 12 October 1840.
			
			``Other clues further narrow the likely time of copying to mid-1839. Young's journal entries between 14 September 1839 and September 1840 are inscribed in blue ink or in darker ink than that used for the 8 July revelation, suggesting he copied the revelation prior to 14 September. Additionally, Young used the volume on 14 July 1839 to record at least one person to contact when he arrived in England to proselytize, suggesting that during summer 1839, Young was already planning to take his journal with him on his mission. At the time, other apostles, such as Wilford Woodruff, were copying JS revelations and discourses into personal volumes that they intended to bring on their mission. It seems plausible that Young also copied relevant texts while preparing for the overseas mission and that he copied the featured revelation between April and mid-September 1839.''
			
		% --
		\paragraph{HDDC}
		% --
		
			HDDC 1536--1543 has a lot more background here, including more journal entries and excerpts from letters.
					
	% -----
	\subsubsection{Versions}
	% -----
		\label{DC118-V}
		
		See the \href{https://drive.google.com/file/d/1goAvNz8jS4R8slYfMG_db55Ty2z_vmgK/view?usp=sharing}{sections comparison} document.
	
		\begin{compactdesc}
			\item[JSP] \hyperref[EMS-1838-JS-8-JUL-1838-A]{8 July 1838}
			\item[BC] ---
			\item[DC 1835] ---
			\item[DC 1844] ---
			\item[TS] ---
			\item[DN] 3:10 (2 April 1853), 37
			\item[MS] 16:12 (25 March 1854), 184
		\end{compactdesc}

	% -----
	\subsubsection{Section 118:1} % *DC118-1 
	% -----
		\label{DC118-1}

		VERILY, thus saith the Lord: Let a conference be held immediately; let the Twelve be organized; and let men be appointed to supply the place of those who are fallen.

		% \begin{framed}
		% \scriptsize
			% \begin{compactdesc}
				% \item[JSP] 
				% % \item[BC] 
				% % \item[]
			% \end{compactdesc}
		% \end{framed}

	% -----
	\subsubsection{Section 118:2} % *DC118-2 
	% -----
		\label{DC118-2}

		Let my servant Thomas remain for a season in the land of Zion, to publish my word.

		% \begin{framed}
		% \scriptsize
			% \begin{compactdesc}
				% \item[JSP] 
				% % \item[BC] 
				% % \item[]
			% \end{compactdesc}
		% \end{framed}

	% -----
	\subsubsection{Section 118:3} % *DC118-3 
	% -----
		\label{DC118-3}

		Let the residue continue to preach from that hour, and if they will do this in all lowliness of heart, in meekness and humility, and long-suffering, I, the Lord, give unto them a promise that I will provide for their families; and an \hyperref[EFFECTUAL]{effectual} door shall be opened for them, from henceforth.

		% \begin{framed}
		% \scriptsize
			% \begin{compactdesc}
				% \item[JSP] 
				% % \item[BC] 
				% % \item[]
			% \end{compactdesc}
		% \end{framed}

	% -----
	\subsubsection{Section 118:4} % *DC118-4 
	% -----
		\label{DC118-4}

		And next spring let them depart to go over the great waters, and there promulgate my gospel, the fulness thereof, and bear record of my name.

		% \begin{framed}
		% \scriptsize
			% \begin{compactdesc}
				% \item[JSP] 
				% % \item[BC] 
				% % \item[]
			% \end{compactdesc}
		% \end{framed}

	% -----
	\subsubsection{Section 118:5} % *DC118-5 
	% -----
		\label{DC118-5}

		Let them take leave of my saints in the city of Far West, on the twenty-sixth day of April next,%
			\footnote{%
				% \label{}%
				So this would be 26 April 1839. See \hyperref[EXSS-DC118-WWJ]{this extended study} for WW's journal entry on that date, which records a council held by the apostles in Far West in preparation for their mission.
				} 
		on the building-spot of my house, saith the Lord.

		% \begin{framed}
		% \scriptsize
			% \begin{compactdesc}
				% \item[JSP] 
				% % \item[BC] 
				% % \item[]
			% \end{compactdesc}
		% \end{framed}

	% -----
	\subsubsection{Section 118:6} % *DC118-6 
	% -----
		\label{DC118-6}

		Let my servant John Taylor, and also my servant John E. Page, and also my servant Wilford Woodruff, and also my servant Willard Richards, be appointed to fill the places of those who have fallen, and be officially notified of their appointment.

		% \begin{framed}
		% \scriptsize
			% \begin{compactdesc}
				% \item[JSP] 
				% % \item[BC] 
				% % \item[]
			% \end{compactdesc}
		% \end{framed}